\section{Limitations, Threats to Validity, and Conclusion}

\subsection{Research Limitations}
\begin{enumerate}
  \item \textbf{Benchmark Simulation}: Evaluations are conducted on synthetic NASA C-MAPSS data where degradation is monotonic. Operational aircraft engines experience maintenance reversibility and non-engine sensor faults not captured in the benchmark.
  \item \textbf{Synthetic Economic Scenarios}: The Simulated Cost Units ($\CPM=1.0, \CFAIL=5.0, \CWASTE=0.20$) represent study-defined benchmark preferences, not empirical airline accounting constants.
  \item \textbf{Conformal Theory Boundaries}: Hierarchical conformal theory provides marginal predictive coverage under its stated sampling and exchangeability assumptions. IntelliTwin uses an engine-weighted construction motivated by that theory. Applicability of those assumptions to ordered, lifetime-stratified degradation trajectories is not established merely by weighting engines equally; reliability is therefore evaluated empirically under the frozen protocol. Furthermore, downstream sequential stopping times are not theoretically guaranteed by split conformal guarantees.
  \item \textbf{Research Prototype Boundaries}: The software is a research prototype evaluating an offline, trace-driven simulation replay engine; it makes no claim of production readiness, aircraft certification, physical hardware integration, or independently established end-to-end operational reproducibility.
\end{enumerate}

\subsection{Conclusion and Next Milestones}
IntelliTwin establishes an end-to-end, scientifically reproducible research architecture linking prognostic neural networks, hierarchical conformal calibration, and sequential maintenance simulation. All methodology, base models, calibrator bundles, and evaluation contracts are completely frozen and audited. Upon completion of the official held-out test evaluation, verified empirical results will be incorporated to resolve RQ1--RQ4.
